\section{Discussion}
\label{sec:discussion}

Our study examines character formation in miniature in the context of emergent misalignment.
In our results, a narrow finetune spills over to other bad behaviors when pretraining has tied behavior to a persistent identity.
Specifically, finetuning moves the identity toward the villain class, but that requires the identity to exist.
This is aligned with the persona-selection explanation of emergent misalignment \citep{marks2026persona}, though a step removed from chat models and agents, where the assistant persona itself is the relevant character.
In our setting, the pretraining corpus determines whether the finetune spreads, a link proposed by \citet{soligo2026emergent} but left open.
There are also alternative explanations for emergent misalignment, which we do not test, such as \citet{li2026emergent}, who predict the spread from the base model's distance to the training data.

Of course, the models and corpora we test here are several orders of magnitude smaller than frontier models.
Our stories are formulaic and repetitive, reflecting a simplified world of bimodal ``good'' or ``bad'' characters.
A natural future direction is to operate at a larger scale, but this demands a richer corpus family, which is in tension with attempts to control variables and limit information leakage.
Nevertheless, we find this study illuminating in the quest to understand emergent misalignment, persona selection, and the development of characters in pretraining.
